\section{Architecture \texttt{uniVALplc}}
\label{secUnivalplcIntro}

La conversion des données du robot (section~\ref{secConversionRobot})
repose sur une architecture logicielle particulière, propre à Stäubli,
qu'il faut rappeler avant d'en détailler les blocs fonctionnels.

\begin{UPSTIinfor}{Architecture client/serveur \texttt{uniVALplc}}
\begin{description}
    \item[Serveur] Installé sur le contrôleur \textbf{CS9} du robot. Il
    répond aux commandes reçues de l'automate M340 via Ethernet/IP, les
    interprète et génère les mouvements correspondants. Il peut faire
    référence à des objets internes du contrôleur CS9 (points, outils,
    repères...).
    \item[Client] Bibliothèque de blocs fonctionnels fournie par le
    fabricant de l'automate (ici Schneider, pour Control Expert). C'est
    cette bibliothèque qui expose les blocs \texttt{VAL\_} et les types
    associés (\texttt{T\_FromRobot}, \texttt{T\_StaeubliRobot}...)
    utilisés dans ce document.
\end{description}
\end{UPSTIinfor}

Pour que cette intégration fonctionne, deux conditions de configuration
doivent être respectées (indépendamment des blocs eux-mêmes) :

\begin{itemize}
    \item l'ordre des octets (\emph{endianness}, section~\ref{secEndianness})
    configuré sur le serveur doit correspondre à celui de l'automate
    client ;
    \item le sélecteur de mode de marche du robot (pupitre \texttt{WMS9})
    doit être positionné sur \textbf{« Commande à distance »}, sans quoi
    les commandes envoyées par l'automate sont ignorées par le serveur.
\end{itemize}

\begin{UPSTIinfor}{Deux familles de blocs pour commander le robot}
Comme rappelé en section~\ref{secRappelBlocsFonctionnels}, les blocs
disponibles pour agir sur le robot se répartissent en deux familles :
\begin{itemize}
    \item les blocs \texttt{VAL\_}, spécifiques à \texttt{uniVALplc}, qui
    donnent accès aux fonctionnalités propres au serveur (lecture/écriture
    de l'état global, accès direct aux E/S de la baie CS9 via
    \texttt{VAL\_ReadWriteIO}, calcul de positions relatives via
    \texttt{VAL\_ShiftPoint}) ;
    \item les blocs \texttt{MC\_}, conformes au standard \textbf{PLCopen},
    qui commandent les mouvements et l'état du groupe d'axes de façon
    normalisée (indépendante du fabricant de robot).
\end{itemize}
Les deux familles partagent le \textbf{même} objet \texttt{AxesGroup}
(\texttt{T\_StaeubliRobot}) : c'est ce partage qui permet à un bloc
\texttt{MC\_} de connaître l'état lu par \texttt{VAL\_ReadAxesGroup} et de
préparer une commande qui sera transmise par \texttt{VAL\_WriteAxesGroup}.
\end{UPSTIinfor}

\begin{UPSTIactivite}[4][Client, serveur et familles de blocs][][][TD]
\UPSTIquestion{Les blocs \texttt{VAL\_} \textbf{et} les blocs \texttt{MC\_}
s'exécutent tous les deux côté \textbf{client}, dans le programme Control
Expert de l'automate M340 : aucun des deux ne \og tourne\fg{} sur la baie
CS9. Expliquer alors, en une phrase, ce que \textbf{le serveur
\texttt{uniVALplc}} (sur la CS9) apporte réellement dans cette
architecture, si ce n'est pas d'exécuter lui-même les blocs.}
\UPSTIquadrillage{3}

\UPSTIquestion{Pour chacune des tâches suivantes, indiquer si elle relève
de la famille \texttt{VAL\_} ou de la famille \texttt{MC\_}, en
justifiant par le \textbf{rôle} de chaque famille (et non en devinant le
nom du bloc) :
\begin{enumerate}
    \item calculer une position décalée par rapport à un point de
    référence enregistré dans le contrôleur CS9 ;
    \item commander un déplacement linéaire absolu du TCP vers une
    position cartésienne donnée ;
    \item accéder directement, depuis l'automate, à une entrée/sortie
    physique de la baie CS9, indépendamment de tout mouvement ;
    \item remettre le groupe d'axes sous puissance après un arrêt.
\end{enumerate}}
\UPSTIquadrillage{4}

\UPSTIquestion{Après une intervention de maintenance, l'automate M340
communique bien avec la baie CS9 (bits de santé du coupleur
\texttt{NOC\_ETHIP\_ROBOT}, section~\ref{secConversionHmi}, corrects),
mais le robot n'exécute aucune des
commandes envoyées. Citer les deux points de configuration, indépendants
du programme automate, à vérifier en priorité, et expliquer pour chacun
la conséquence d'un mauvais réglage.}
\UPSTIquadrillage{4}
\end{UPSTIactivite}
